\chapter{Declaration on the use of generative AI and TRACE}
\label{app:ai}

This appendix is the declaration I will defend. It is also the TRACE of the AI Studio origin of the first application.

\section{Declaration}
\label{sec:ai-decl}

I used generative AI in this project. I did not hide it.

\begin{itemize}
\item Google AI Studio generated the first web application (React tabs, the first Express routes, and the first visual layout).
\item An earlier Studio version also included an academic-report Word export. Those leftover \texttt{.docx} files were removed from this repository. I did not keep a generator that writes the official report.
\item Gemini is still used on the teaching email routes \texttt{/api/generate-phishing} and \texttt{/api/analyze-email}, and on non-preset deepfake filenames. Each of those answers carries \texttt{source} equal to \texttt{gemini}, \texttt{fallback}, or \texttt{preset}.
\item The adversarial route does not use Gemini. When the Python service is down, the answer is 503, not an invented confidence.
\item I implemented the lock, the tests, and the evaluation wiring in this repository.
\item This report text is written in the official UPB \texttt{project-report} template. I have to read every sentence and every citation before the viva. Every number must be checked by me, Rinor Basholli.
\end{itemize}

The work could not exist in this form without my decisions: the requirements, the source labels, the refusal to harvest mail, the CIFAR model choice, the unused adversarial fallback, and the honesty about the first F1 $= 1.0$ run and the later 0.756 wording run. The same work could have been written without AI, more slowly, as ordinary React and Python. AI did not sit the course. I did.

\section{TRACE of Studio metadata}
\label{sec:trace}

The repository still shows that it began in AI Studio.

\paragraph{metadata.json}
The file \texttt{metadata.json} at the project root is Studio project metadata. The name is ``AI Cybersecurity Risk Lab''. The description talks about an interactive study studio. \texttt{majorCapabilities} contains \texttt{MAJOR\_CAPABILITY\_SERVER\_SIDE\_GEMINI\_API}. \texttt{requestFramePermissions} is an empty list. I left this file in place so that the origin is visible.

\paragraph{vite.config.ts}
The Vite config still has comments about HMR being disabled in AI Studio through \texttt{DISABLE\_HMR}. That flag exists so that preview edits do not flicker. I did not invent it.

\paragraph{aistudio-build}
AI Studio publishes a Node website. The deploy notes in the project README warn that a republish from Studio can overwrite a custom all-in-one image. There is no separate \texttt{aistudio-build/} source tree in the current repository. The remaining traces are \texttt{metadata.json}, the HMR comments, the Gemini-first routes, and the leftover UI labels (the ``DB ACTIVE'' footer and the ResNet-50 button).

\paragraph{Removed Word export}
The academic-report \texttt{.docx} path is gone. I mention the removal here so that an examiner who saw an older tree knows I deleted it on purpose. The official report is this \LaTeX\ tree under \texttt{report/}, using BibTeX \texttt{ieeetr} as the template requires.

\section{Model identifier}
\label{sec:model-id}

The default Gemini model id in \texttt{server.ts} is \texttt{gemini-3.5-flash}. The environment variable \texttt{GEMINI\_MODEL} overrides it. I must say in the viva which id the live key served. I do not invent a transcript of those calls in this appendix.
